\scanpage{26}

\diagram{\begin{tikzcd}[ampersand replacement=\&,column sep=small,row sep=small]
\operatorname{Hom}(\Phi(i,n),\Psi^j)\arrow[ddd] \& \subset\operatorname{Hom}(\Phi(i,n),\Phi^{j+1})\\
\& \begin{matrix}\Vert\\\Phi(j+1,\unsure{n|1|-})\end{matrix}\\
=0\ \text{si} \& j+1\ne i,i+1\quad\text{i.e. si }j\ne i-1,i\\
\operatorname{Hom}(\Psi^{i-1},\Psi^j)=0\ \text{si }j\ne i-1 \&
\end{tikzcd}}

Si donc $j=i$, $\operatorname{Hom}(\Phi(i,n),\Psi^{j=i})\xleftarrow{\sim}\operatorname{Hom}(\unsure{\Psi^i|\Phi^i},\Psi^i)$ $\underset{k}{\unsure{\simeq|=}}$ (base l'hom can)

\deleted{Si} Si $j=i-1$,

\diagram{\begin{tikzcd}[ampersand replacement=\&,column sep=small,row sep=small]
{}-{}\& \Psi^{i-1}\arrow[r]\arrow[d,"n\,\mathrm{id}"'] \& \Phi^i\deleted{\unsure{(1)|(k)}}\arrow[r]\arrow[d] \& \Psi^i\arrow[r,dash] \& 0\\
0\arrow[r,dash] \& \Psi^{i-1}\arrow[r,"\unsure{\iota|i|v}"] \& \Phi(i,n)\arrow[r,dash]\arrow[d,"\alpha=?"] \& \Psi^i\arrow[r,dash]\arrow[u] \& 0\\
\& \& \Psi^{i-1} \& \&
\end{tikzcd}}

le composé de $\alpha\deleted{\unsure{\iota|i}}$ avec \deleted{$\unsure{n\,\mathrm{id}_{\Psi^i}|u\circ d_{\Psi^i}}$} \inserted{$\Phi^i\mathrel{-}\Phi(i,n)$} est $\in\operatorname{Hom}(\Phi^i,\Psi^{i-1})=0$, donc \deleted{$n\alpha\unsure{\iota|i}\in\operatorname{Hom}$} \inserted{son composé} $\alpha\iota\circ n\,\mathrm{id}_{\Psi^{i-1}}=n(\alpha\iota)\in\operatorname{Hom}(\Psi^{i-1},\Psi^{i-1})$ est nul, et comme $\operatorname{Hom}(\Psi^{i-1},\Psi^{i-1})\simeq k$, on voit que $\alpha\iota=0$, donc $\alpha$ \deleted{\unsure{passe par|provient}} se factorise par $\Psi^i\mathrel{-}\Psi^{i-1}$, donc est nul. Donc

\[\operatorname{Hom}(\Phi(i,n),\Psi^j)=
\begin{cases}0&\text{si }j\ne i\\
\simeq k\ (\text{base l'hom can }\Phi(i,n)\mathrel{-}\Psi^i)&\text{si }j=i\end{cases}\]

% archive sheet number: 25
